\chapter{Uvod}
\label{ch:uvod}

Razvoj novih lijekova jedan je od najskupljih i vremenski najzahtjevnijih
procesa u suvremenoj znanosti. Klasičan pristup, u kojem se kandidati za
djelatne tvari pronalaze laboratorijskim pretraživanjem velikog broja
spojeva, sve se češće nadopunjuje računalnim metodama koje prostor mogućih
molekula suzuju prije nego što se uopće pristupi eksperimentalnoj sintezi.
Među takvim metodama posebno se ističu generativni modeli dubokog učenja,
koji ne biraju samo između postojećih spojeva, već \emph{de novo} predlažu
potpuno nove molekulske strukture s traženim svojstvima.

Posebno zanimljivu skupinu ciljnih molekula čine antimikrobni peptidi
(eng. \emph{antimicrobial peptides}, AMP). To su u vodi topljivi peptidi
sastavljeni od kratkih lanaca aminokiselina, koje višestanični organizmi
prirodno proizvode kao prvu liniju obrane od virusa, bakterija, gljivica i
tumorskih stanica \cite{negovetic2022}. Budući da su uglavnom pozitivno
nabijeni, AMP-ovi djeluju izravno na membrane ciljanih mikroorganizama, a
zanimanje za njih dodatno raste jer dosad nije zabilježena bakterija
otporna na ljudske AMP-ove, što ih čini obećavajućom alternativom
klasičnim antibioticima na koje se otpornost stalno povećava. Upravo zbog
toga generiranje novih peptidnih struktura s visokom predviđenom
antimikrobnom aktivnošću predstavlja vrijedan i aktualan istraživački cilj.

Polazna točka ovog rada je model REINVENT, generativni okvir temeljen na
povratnoj neuronskoj mreži (eng. \emph{recurrent neural network}, RNN) i
pojačanom učenju (eng. \emph{reinforcement learning}, RL). REINVENT najprije
uči distribuciju valjanih molekula iz velike baze spojeva, a zatim se
dobiveni model (tzv. \emph{prior}) usmjerava prema željenom kemijskom
prostoru tako da se nagrađuju molekule koje postižu visok rezultat zadane
funkcije vrednovanja (eng. \emph{scoring function}). Model je, međutim,
izvorno osmišljen i treniran za područje malih organskih molekula nalik
lijekovima, a ne za peptide. Cilj je ovog rada istražiti prednosti i
nedostatke takvog modela te ispitati može li se njegova arhitektura
prilagoditi domeni aminokiselina, odnosno usmjeriti generiranje prema
valjanim peptidnim strukturama zapisanim u formatu SMILES
(eng. \emph{Simplified Molecular Input Line Entry System}).

Da bi se generativni model uopće mogao usmjeriti prema antimikrobnim
peptidima, potrebna je funkcija vrednovanja koja prepoznaje upravo to
svojstvo. Standardne funkcije vrednovanja ugrađene u REINVENT nisu
prikladne jer su vezane uz sličnost s referentnim malim molekulama ili uz
svojstva nevezana za peptide. Stoga je kao funkcija vrednovanja iskorišten
postojeći model strojnog učenja razvijen u srodnom radu \cite{negovetic2022},
treniran upravo na aminokiselinskim sekvencama radi predviđanja
antimikrobne aktivnosti. Time je dobivena prirodna poveznica između dvaju
modela: REINVENT predlaže nove peptidne kandidate, dok ih spomenuti model
ocjenjuje s obzirom na vjerojatnu antimikrobnu aktivnost.

Navedeni referentni model temelji se na algoritmu slučajnih šuma
(eng. \emph{random forest}), ansamblu velikog broja stabala odluke koji
dobro podnosi šum i neuravnotežene skupove podataka te jednako kvalitetno
rješava klasifikacijske i regresijske probleme \cite{negovetic2022}. Model
ne radi izravno nad tekstualnim zapisom molekule, već nad skupom
brojčanih molekulskih značajki (deskriptora) izračunatih iz SMILES zapisa.
Konkretno, koristi se klasifikator s 600 stabala, kod kojeg je najveći broj
značajki po čvoru ograničen na korijen ukupnog broja značajki, a vrednovanje
je provedeno unakrsnom validacijom uz metriku ROC-AUC (eng. \emph{Receiver
Operating Characteristic -- Area Under Curve}). Kako bi broj ulaznih
značajki ostao obradiv, u izvornom je radu nad njima provedena predobrada i
odabir reprezentativnog podskupa. U ovom radu taj se model koristi kao
gotova, prethodno istrenirana komponenta za vrednovanje generiranih peptida.

Pri odabiru konkretne inačice REINVENT okvira odlučeno je raditi s
ranijom, arhitektonski jednostavnijom verzijom modela umjesto s
najsuvremenijom. Razlog je tome jasnija i preglednija struktura programskog
koda s manje vanjskih ovisnosti, što je olakšalo razumijevanje i izmjenu
same petlje pojačanog učenja, kao i ugradnju vlastite funkcije vrednovanja.
Naprednije inačice nude dodatne načine rada (primjerice dekoraciju okosnice
molekule ili povezivanje fragmenata) koji za ovaj zadatak nisu nužni, pa bi
njihova složenost unijela nepotrebne troškove bez stvarne koristi.
Mogućnost korištenja modela strojnog učenja izvan izvorne arhitekture
REINVENT-a time je iskorištena u dva smjera: zamjenom ugrađene funkcije
vrednovanja vanjskim modelom slučajnih šuma te uvođenjem dodatnih
ograničenja kojima se generiranje usmjerava prema valjanim peptidima.

Sukladno opisanom, ovaj rad nastoji odgovoriti na sljedeće glavno
istraživačko pitanje: \emph{može li se generativni model REINVENT, izvorno
namijenjen malim organskim molekulama, prilagoditi tako da pouzdano generira
valjane peptidne strukture s visokom predviđenom antimikrobnom aktivnošću,
te uz koje je intervencije u funkciju vrednovanja to moguće postići?} Iz tog
glavnog pitanja proizlazi nekoliko užih pod-pitanja na koja se odgovara kroz
provedene eksperimente:

\begin{itemize}
  \item Je li za usmjeravanje generiranja prema peptidima dovoljno samo
        zamijeniti funkciju vrednovanja, ili je nužno i eksplicitno
        filtriranje kojim se nepeptidne molekule odbacuju?
  \item Može li se prethodno istreniran model strojnog učenja, razvijen
        izvan izvorne arhitekture REINVENT-a, uspješno uklopiti kao vanjska
        funkcija vrednovanja?
  \item Kako strogost dodatnih ograničenja (najmanja duljina peptida,
        najveći udio nekanonskih aminokiselina te zabrana ponavljanja istih
        aminokiselina) utječe na sastav i kvalitetu generiranih peptida?
  \item Daje li bolje rezultate tvrdo odbacivanje ili blaga kazna
        nekanonskih aminokiselina u funkciji vrednovanja?
  \item Nasljeđuje li agent inicijaliziran iz prethodno istreniranog modela
        (tzv. topli start) pristranosti tog modela?
  \item Koliko je proces učenja pojačanim učenjem ponovljiv s obzirom na
        njegovu stohastičku narav?
\end{itemize}

U skladu s tim pitanjima, u radu je ispitana mogućnost prenamjene postojećeg
\emph{prior} modela za generiranje aminokiselinskih sekvenci, razvijeno je
nekoliko varijanti dodatnih filtara i kazni kojima se oblikuje funkcija
vrednovanja, te su provedeni i međusobno uspoređeni eksperimenti treniranja
agenta. Dobiveni rezultati vrednovani su s obzirom na udio stvarno
generiranih peptida, njihov sastav te kretanje rezultata vrednovanja kroz
proces učenja, a potom su dovedeni u vezu s relevantnim rezultatima iz
srodne literature.
